\section{仿真验证}

\subsection{实验设置}

\textbf{仿真平台}：基于MATLAB的XHY AUV 6自由度动力学仿真器，使用CFD RANS标定的阻力模型（表\ref{tab:cfd}）作为名义动力学，时间步长$\Delta t=0.01$s，仿真时长$T=100$s。

\textbf{海流场景}：
\begin{enumerate}
    \item \textbf{圆形轨迹}（持续机动）：AUV以$u_d=1$ m/s跟踪半径50m的圆形参考轨迹，持续激励偏航机动。
    \item \textbf{直线轨迹}（低激励）：AUV以$u_d=1$ m/s沿直线参考轨迹航行，验证激励门控在弱激励下的表现。
    \item \textbf{海流阶跃}（突变）：$t=50$s时海流从$V_c=0.15$ m/s、$\beta_c=30^\circ$阶跃至$V_c=0.45$ m/s、$\beta_c=60^\circ$，验证观测器对海流突变的响应速度。
\end{enumerate}

所有场景中海流由一阶Gauss-Markov过程生成（相关时间常数$\tau_c=100$s，均值$V_c=0.3$ m/s，$\beta_c=45^\circ$，标准差$\sigma_V=0.1$ m/s）。

\textbf{模型失配}：通过在Plant模型中随机扰动CFD阻力系数来模拟模型失配。每个自由度独立扰动，扰动幅度为0\%、10\%、20\%、30\%。观测器始终使用名义CFD模型。

\textbf{对比方法}：
\begin{itemize}
    \item \textbf{KIN}：运动学海流观测器\cite{liang2018three}，作为不依赖动力学模型的基线
    \item \textbf{LESO}：标准线性扩张状态观测器，估计集总扰动
    \item \textbf{PIESO}：物理信息ESO，嵌入Gauss-Markov海流先验
    \item \textbf{EKF-tuned}：CFD增广扩展卡尔曼滤波（14状态：$\bm{\nu}+\bm{c}+$力偏置），过程噪声协方差Q按失配水平自适应调优
    \item \textbf{UCCO}：本文提出的EG-UCCO（仅Yaw前馈）
\end{itemize}

\textbf{评估指标}：海流估计的均方根误差RMSE$_{V_c}=\sqrt{\frac{1}{N}\sum_{i=1}^{N}\|\hat{\bm{c}}_i-\bm{c}_i^{\text{true}}\|^2}$（m/s），平均绝对误差MAE$_{V_c}$。

\subsection{海流估计精度}

表\ref{tab:accuracy}给出了3种场景、4种模型失配水平下各方法的RMSE$_{V_c}$（2个随机种子的均值）。

\begin{table}[h]
\centering
\caption{海流估计精度对比 (RMSE$_{V_c}$, m/s)}
\label{tab:accuracy}
\begin{tabular}{llcccc}
\toprule
场景 & 失配 & KIN & EKF-tuned & LESO* & UCCO \\
\midrule
\multirow{4}{*}{圆形} & 0\%   & 0.265 & 0.268 & 0.261 & \textbf{0.142} \\
                      & 10\%  & 0.265 & 0.277 & 0.261 & \textbf{0.149} \\
                      & 20\%  & 0.266 & 0.285 & 0.261 & \textbf{0.144} \\
                      & 30\%  & 0.265 & 0.290 & 0.261 & \textbf{0.153} \\
\midrule
\multirow{4}{*}{直线} & 0\%   & 0.291 & 0.284 & 0.280 & \textbf{0.166} \\
                      & 10\%  & 0.291 & 0.281 & 0.280 & \textbf{0.171} \\
                      & 20\%  & 0.291 & 0.285 & 0.280 & \textbf{0.171} \\
                      & 30\%  & 0.291 & 0.286 & 0.280 & \textbf{0.152} \\
\midrule
\multirow{4}{*}{阶跃} & 0\%   & 0.336 & 0.204 & 0.335 & \textbf{0.200} \\
                      & 10\%  & 0.336 & 0.215 & 0.335 & \textbf{0.198} \\
                      & 20\%  & 0.336 & 0.227 & 0.335 & \textbf{0.198} \\
                      & 30\%  & 0.336 & 0.237 & 0.335 & \textbf{0.197} \\
\bottomrule
\end{tabular}
\\[4pt]
\footnotesize{*LESO/PIESO估计集总扰动而非物理海流，其RMSE$_{V_c}\approx\bar{V}_c$反映的是海流均值基线。}
\end{table}

\textbf{核心发现}：
\begin{enumerate}
    \item UCCO在全部12个场景-失配组合中估计精度最优（粗体标注）。
    \item 在圆形机动场景（持续激励），UCCO的RMSE$_{V_c}$约为0.14--0.15 m/s，比KIN运动学观测器（0.27 m/s）和调优EKF（0.27--0.29 m/s）提升约45\%。
    \item 在直线低激励场景，UCCO的精度优势依然保持（0.17 vs 0.28--0.29 m/s），说明激励门控在弱激励下不会因停止更新而显著劣化。此时UCCO主要依赖GM时间传播维持估计。
    \item 在海流阶跃场景，UCCO与EKF-tuned性能接近（0.20 vs 0.20--0.24），均能快速响应海流突变。
\end{enumerate}

\subsection{模型失配鲁棒性}

图\ref{fig:degradation}展示了各方法在圆形轨迹下RMSE$_{V_c}$随模型失配水平的退化趋势。

\begin{figure}[h]
\centering
\fbox{\parbox{0.85\textwidth}{
\textbf{模型失配退化曲线（圆形轨迹）:}

\begin{tabular}{lcccc|c}
\toprule
方法 & 0\% & 10\% & 20\% & 30\% & 退化幅度 \\
\midrule
UCCO      & 0.200 & 0.207 & 0.207 & 0.208 & \textbf{+4.1\%} \\
EKF-tuned & 0.284 & 0.285 & 0.290 & 0.297 & +4.5\% \\
KIN        & 0.265 & 0.265 & 0.266 & 0.265 & 0.0\% \\
LESO/PIESO & 0.334 & 0.334 & 0.334 & 0.334 & 0.0\% \\
\bottomrule
\end{tabular}
}}
\caption{模型失配退化对比}
\label{fig:degradation}
\end{figure}

UCCO的退化幅度（+4.1\%）与调优EKF（+4.5\%）相当，但UCCO的绝对精度始终高于EKF-tuned约30\%。KIN和LESO因不使用CFD模型而天然免疫模型失配，但其绝对精度较低。值得指出的是，未调优的EKF（标称Q矩阵）在相同条件下RMSE$_{V_c}$退化高达694\%，从0.062崩溃至0.492。这从反面验证了参数调优对EKF类方法的敏感依赖性，以及UCCO梯度更新机制的固有鲁棒性。

\subsection{门控消融实验}

为验证激励门控的实际作用，我们在3种传感器噪声水平下对比了3种门控配置。结果如表\ref{tab:gate}所示。

\begin{table}[h]
\centering
\caption{激励门控消融实验 (RMSE$_{V_c}$, m/s, 圆形轨迹)}
\label{tab:gate}
\begin{tabular}{lccc}
\toprule
门控配置 & 无噪声 & 低噪声 & 高噪声 \\
\midrule
Gramian门控 ($\mu=10^{-8}$) & \textbf{0.195} & \textbf{0.263} & \textbf{0.512} \\
无门控 (持续更新)           & 0.202 & 0.278 & 0.668 \\
高阈值 ($\mu=0.01$)          & 0.334 & 0.334 & 0.334 \\
\midrule
门控收益 (vs 无门控)         & +3.6\% & +5.6\% & \textbf{+23.3\%} \\
\bottomrule
\end{tabular}
\end{table}

\textbf{关键观察}：
\begin{enumerate}
    \item Gramian门控在所有噪声水平下均优于无门控方案，且优势随噪声增大而增加：无噪声时+3.6\%，低噪声时+5.6\%，高噪声时+23.3\%。这印证了门控机制在高噪声、弱激励条件下的保护作用。
    \item 高阈值门控($\mu=0.01$)导致激励率降至0\%，海流估计完全被阻断，RMSE$_{V_c}=0.334$恰好等于海流均值基线。这验证了门控阈值选择的重要性。
    \item Gramian门控的激励率为100\%（圆形轨迹持续机动），但门控仍在每次更新中通过Gramian条件数自适应调整有效增益（高条件数$\rightarrow$高正则化$\rightarrow$小步长）。
\end{enumerate}

\subsection{传感器噪声鲁棒性}

在低传感器噪声条件下（DVL速度噪声$\sigma_v=0.02$ m/s，航向角噪声$\sigma_\psi=0.5^\circ$），使用10个随机种子进行了Monte Carlo验证。统计结果见表\ref{tab:montecarlo}。

\begin{table}[h]
\centering
\caption{Monte Carlo验证 (10 seeds, Mean$\pm$Std, RMSE$_{V_c}$ m/s)}
\label{tab:montecarlo}
\begin{tabular}{llccc}
\toprule
场景 & 失配 & KIN & EKF-tuned & UCCO \\
\midrule
\multirow{2}{*}{圆形} & 0\%   & 0.272$\pm$0.062 & 0.276$\pm$0.028 & 0.279$\pm$0.056 \\
                      & 20\%  & 0.272$\pm$0.062 & 0.292$\pm$0.017 & \textbf{0.270$\pm$0.048} \\
\midrule
\multirow{2}{*}{直线} & 0\%   & 0.278$\pm$0.063 & 0.282$\pm$0.017 & 0.306$\pm$0.057 \\
                      & 20\%  & 0.280$\pm$0.063 & 0.290$\pm$0.015 & \textbf{0.276$\pm$0.061} \\
\bottomrule
\end{tabular}
\end{table}

在传感器噪声下，无模型失配(0\%)时各方法性能接近（差异在1--8\%以内）。然而，当模型失配增加到20\%时，UCCO在圆形和直线场景中均取得最优估计（0.270和0.276 m/s）。更重要的是，UCCO从0\%到20\%失配的性能变化为：圆形-3.5\%（改善）、直线-9.8\%（改善），而EKF-tuned分别退化了+5.8\%和+3.2\%。这组结果强有力地支撑了UCCO的核心主张：\textbf{在模型失配条件下提供鲁棒的海流估计}。

需要诚实指出的是，UCCO的估计方差(std$\approx$0.05--0.06)略高于EKF-tuned(std$\approx$0.02--0.03)，表明UCCO的梯度更新机制在噪声下具有较大的逐次波动。这是梯度方法的固有特性，可通过增加滑动窗口平滑或降低自适应增益来缓解。
